\documentclass[11pt]{article}
\usepackage[utf8]{inputenc}
\usepackage[T1]{fontenc}
\usepackage[margin=1in]{geometry}
\usepackage{amsmath,amssymb,mathtools}
\usepackage[ruled,vlined,linesnumbered]{algorithm2e}
\usepackage{enumitem}

\begin{document}

\section*{Integrated Algorithm for the Full Constrained Portfolio Problem}

\subsection*{1. Problem solved by the code}

The implementation solves the reduced portfolio problem
\[
\min_{x \in \mathbb{R}^n}
\left\{
\sigma^2 x^\top \Sigma x - \mu^\top x + \frac{1}{\gamma} g(x)
\right\}
\quad
\text{s.t.}
\quad
x \in \Delta_n,
\qquad
l \le Ax \le u,
\]
where
\[
\Delta_n := \{x \in \mathbb{R}^n : x \ge 0,\ \mathbf{1}^\top x = 1\}.
\]

The penalty \(g(x)\) is the reduced form of the lifted model with auxiliary
variables \(z\):
\[
g(x)
=
\frac12 \min_{z}
\left\{
\sum_{i=1}^n \frac{x_i^2}{z_i}
:
z_i \in [0,1],\quad
\sum_{i=1}^n z_i \le k,\quad
z_i \ge x_i
\right\}.
\]
Intuitively:
\begin{itemize}[topsep=0.3em,itemsep=0.2em]
\item \(x \in \Delta_n\) keeps the portfolio long-only and fully invested;
\item \(l \le Ax \le u\) imposes extra exposure bounds;
\item \(g(x)\) is the nonlinear portfolio regularizer controlled by \(k\);
\item \(\gamma\) scales how strongly that regularizer is enforced.
\end{itemize}

The code handles the linear exposure constraints by an outer augmented
Lagrangian loop, and handles the simplex-constrained proximal step by a custom
inner routine.

\subsection*{2. Constraint encoding}

The two-sided exposure constraints are rewritten as
\[
B :=
\begin{bmatrix}
A \\
-A
\end{bmatrix},
\qquad
b :=
\begin{bmatrix}
u \\
-l
\end{bmatrix},
\qquad
Bx - b \le 0.
\]

\subsection*{3. Outer proximal augmented Lagrangian model}

At outer iteration \(r\), the code forms the smooth model
\[
\mathcal{L}_r(x)
:=
\sigma^2 x^\top \Sigma x - \mu^\top x
+
\frac{1}{2\rho_r}\left\|[p^r + \rho_r(Bx-b)]_+\right\|_2^2
-
\frac{1}{2\rho_r}\|p^r\|_2^2
+
\frac{1}{2\eta_r}\|x-x^r\|_2^2,
\]
and leaves the nonsmooth/simplex part in
\[
f_2(x) := \frac{1}{\gamma} g(x) + \delta_{\Delta_n}(x).
\]

Hence the primal subproblem is
\[
x^{r+1} \approx \arg\min_x \left\{ \mathcal{L}_r(x) + f_2(x) \right\}.
\]

The implementation solves this subproblem by FISTA. The key point is that the
proximal step is \emph{not} a plain simplex projection; it is the prox of
\(\tau g + \delta_{\Delta_n}\), evaluated by Moreau decomposition, PAVA, and a
breakpoint walk in the simplex multiplier.

\subsection*{4. How to read the inner prox}

Given a gradient step \(v\), the inner routine must compute
\[
\operatorname{prox}_{\tau g + \delta_{\Delta_n}}(v)
=
\arg\min_{x \in \Delta_n}
\left\{
\frac12 \|x-v\|_2^2 + \tau g(x)
\right\}.
\]

This is handled in two layers.

\paragraph*{Free prox without the simplex equality.}
First imagine dropping the constraint \(\mathbf{1}^\top x = 1\). Then we only
need \(\operatorname{prox}_{\tau g}(w)\), which the code computes by Moreau's
identity:
\[
\operatorname{prox}_{\tau g}(w)
=
w - \tau \operatorname{prox}_{g^\ast/\tau}(w/\tau).
\]
The conjugate-side prox \(\operatorname{prox}_{g^\ast/\tau}\) has a sorted
block structure and is evaluated efficiently by PAVA.

\paragraph*{Restoring the simplex equality.}
To impose \(\mathbf{1}^\top x = 1\), introduce a scalar multiplier \(\nu\) and
define
\[
\Phi(\nu) := \mathbf{1}^\top \operatorname{prox}_{\tau g}(v - \nu \mathbf{1}).
\]
The required prox is obtained by finding the unique root of
\[
\Phi(\nu) = 1.
\]

The expensive part would be recomputing the free prox from scratch for many
trial values of \(\nu\). The code avoids this by observing that, after sorting,
the path of \(\operatorname{prox}_{\tau g}(v - \nu \mathbf{1})\) changes only
at finitely many breakpoints. Between breakpoints, \(\Phi(\nu)\) is affine.

So the algorithm:
\begin{itemize}[topsep=0.3em,itemsep=0.2em]
\item starts from a right bracket \(\nu_R\) with \(\Phi(\nu_R) \le 1\);
\item moves only toward smaller \(\nu\);
\item tracks the current PAVA block partition;
\item updates only two kinds of events: regime changes of a block, and merges
of adjacent blocks;
\item solves for the exact root as soon as it lies on the current affine
  segment.
\end{itemize}

This is what the function
\texttt{prox\_moreau\_simplex\_breakpoint\_walk} implements.

\begin{algorithm}[H]
\caption{Outer proximal augmented Lagrangian solver}
\label{alg:outer_palm_portfolio}
\KwIn{Data \(\Sigma,\mu,A,l,u,k,\sigma,\gamma\); feasible \(x^0 \in \Delta_n\) with \(l \le Ax^0 \le u\); penalty \(\rho_0 > 0\); proximal parameter \(\eta_0 > 0\); tolerances \(\varepsilon_{\mathrm{outer}}, \varepsilon_{\mathrm{inner}}\)}
\KwOut{Approximate solution \(x^\star\)}

\BlankLine
\tcp{Encode the two-sided exposure constraints}
\(
B \leftarrow
\begin{bmatrix}
A \\
-A
\end{bmatrix}
\),
\(
b \leftarrow
\begin{bmatrix}
u \\
-l
\end{bmatrix}
\)
\;

\BlankLine
\tcp{Initialization}
\(r \leftarrow 0\) \;
\(x^r \leftarrow x^0\) \;
\(p^r \leftarrow 0\) \;
\(\rho_r \leftarrow \rho_0\), \(\eta_r \leftarrow \eta_0\) \;
\(\nu^{\mathrm{warm}} \leftarrow \texttt{None}\) \;
\(E^r \leftarrow \min\{b - Bx^r,\ p^r\}\) \;
\(x^{\mathrm{init}} \leftarrow x^0\) \;

\BlankLine
\While{not converged}{
    \tcp{Current Lipschitz constant used by the inner FISTA loop}
    \[
    L_r \leftarrow 2 \sigma^2 \lambda_{\max}(\Sigma) + \rho_r \|B\|_2^2 + \eta_r^{-1}.
    \]

    \BlankLine
    \tcp{Approximately solve the primal subproblem}
    \(
    (x^{r+1}, \texttt{diag})
    \leftarrow
    \textsc{InnerFISTA}(x^r, p^r, \rho_r, \eta_r, \nu^{\mathrm{warm}}, \varepsilon_{\mathrm{inner}})
    \)
    \;
    \If{\(\texttt{diag.last\_nu}\) is finite}{
        \(\nu^{\mathrm{warm}} \leftarrow \texttt{diag.last\_nu}\) \;
    }

    \BlankLine
    \tcp{Multiplier update for \(Bx-b \le 0\)}
    \[
    p^{r+1} \leftarrow [p^r + \rho_r(Bx^{r+1} - b)]_+ .
    \]

    \BlankLine
    \tcp{Residuals checked by the implementation}
    \[
    r_{\mathrm{stat}} \leftarrow \|x^{r+1} - x^r\|_2,
    \qquad
    r_{\Delta} \leftarrow |\mathbf{1}^\top x^{r+1} - 1|,
    \]
    \[
    r_{\mathrm{exp}} \leftarrow
    \max\!\left\{
    \|(Ax^{r+1} - u)_+\|_\infty,\,
    \|(l - Ax^{r+1})_+\|_\infty
    \right\}.
    \]

    \BlankLine
    \If{\(r_{\mathrm{stat}} \le \varepsilon_{\mathrm{outer}}\) and \(r_{\Delta} \le \varepsilon_{\mathrm{outer}}\) and \(r_{\mathrm{exp}} \le \varepsilon_{\mathrm{outer}}\)}{
        \Return \(x^{r+1}\) \;
    }

    \BlankLine
    \tcp{Residual surrogate used for penalty adaptation}
    \[
    E^{r+1} \leftarrow \min\{b - Bx^{r+1},\ p^r/\rho_r\}.
    \]
    \[
    g_r \leftarrow \rho_0 (r+2)^{\alpha_{\mathrm{growth}}}.
    \]
    \If{\(\|E^r\|_\infty \le 10^{-12}\)}{
        \If{\(\|E^{r+1}\|_\infty > 10^{-12}\)}{
            \(\rho_{r+1} \leftarrow \max\{\xi_\rho \rho_r,\ g_r\}\) \;
        }
        \Else{
            \(\rho_{r+1} \leftarrow \rho_r\) \;
        }
    }
    \Else{
        \If{\(\|E^{r+1}\|_\infty > \beta \|E^r\|_\infty\)}{
            \(\rho_{r+1} \leftarrow \max\{\xi_\rho \rho_r,\ g_r\}\) \;
        }
        \Else{
            \(\rho_{r+1} \leftarrow \rho_r\) \;
        }
    }

    \BlankLine
    \tcp{Proximal regularization update}
    \[
    \eta_{r+1}
    \leftarrow
    \max\!\left\{
    \delta_\eta \|x^{r+1} - x^{\mathrm{init}}\|_2^2,\,
    \eta_0 (r+2)^{\alpha_{\mathrm{growth}}}
    \right\}.
    \]

    \BlankLine
    \(E^r \leftarrow E^{r+1}\) \;
    \(p^r \leftarrow p^{r+1}\) \;
    \(x^r \leftarrow x^{r+1}\) \;
    \(r \leftarrow r + 1\) \;
}

\BlankLine
\Return \(x^r\) \;
\end{algorithm}

\begin{algorithm}[H]
\caption{\(\textsc{InnerFISTA}\): approximate primal subproblem solve}
\label{alg:inner_fista_portfolio}
\KwIn{Center \(x^r\); multipliers \(p^r\); penalty \(\rho_r\); proximal parameter \(\eta_r\); warm start \(\nu^{\mathrm{warm}}\); tolerance \(\varepsilon_{\mathrm{inner}}\)}
\KwOut{Approximate minimizer \(x^{r+1}\); diagnostics}

\BlankLine
\tcp{Initialization}
\[
L_r \leftarrow 2 \sigma^2 \lambda_{\max}(\Sigma) + \rho_r \|B\|_2^2 + \eta_r^{-1},
\qquad
\tau_r \leftarrow \frac{1}{L_r \gamma}.
\]
\[
x^{(0)} \leftarrow x^r,
\qquad
y^{(0)} \leftarrow x^r,
\qquad
a_0 \leftarrow 1.
\]

\BlankLine
\For{\(t=0,1,2,\dots\)}{
    \tcp{Gradient of the smooth model}
    \[
    \nabla \mathcal{L}_r(y^{(t)})
    =
    2 \sigma^2 \Sigma y^{(t)} - \mu
    + B^\top [p^r + \rho_r(By^{(t)} - b)]_+
    + \eta_r^{-1}(y^{(t)} - x^r).
    \]

    \BlankLine
    \tcp{Gradient step}
    \[
    v^{(t)} \leftarrow y^{(t)} - \frac{1}{L_r}\nabla \mathcal{L}_r(y^{(t)}).
    \]

    \BlankLine
    \tcp{Simplex-constrained proximal step}
    \[
    x^{(t+1)}
    \leftarrow
    \operatorname{prox}_{\tau_r g + \delta_{\Delta_n}}(v^{(t)}).
    \]
    Compute this prox by
    \[
    (x^{(t+1)}, \texttt{prox\_diag})
    \leftarrow
    \textsc{SimplexProx}(v^{(t)}, \tau_r, k, \nu^{\mathrm{warm}}).
    \]
    Update \(\nu^{\mathrm{warm}}\) from \(\texttt{prox\_diag.nu}\) \;

    \BlankLine
    \tcp{Stopping rule used in the code}
    \[
    r_{\mathrm{inner}} \leftarrow \|x^{(t+1)} - x^{(t)}\|_2.
    \]
    \If{\(r_{\mathrm{inner}} \le \varepsilon_{\mathrm{inner}}\)}{
        \Return \(x^{(t+1)}\) and diagnostics \;
    }

    \BlankLine
    \tcp{FISTA momentum}
    \[
    a_{t+1} \leftarrow \frac{1 + \sqrt{1 + 4 a_t^2}}{2},
    \qquad
    y^{(t+1)}
    \leftarrow
    x^{(t+1)}
    +
    \frac{a_t - 1}{a_{t+1}}(x^{(t+1)} - x^{(t)}).
    \]
}
\end{algorithm}

\begin{algorithm}[H]
\caption{\(\textsc{SimplexProx}\): Moreau \(+\) PAVA \(+\) directional breakpoint walk}
\label{alg:simplex_prox_portfolio}
\KwIn{Vector \(v\); parameter \(\tau > 0\); budget parameter \(k\); warm start \(\nu^{\mathrm{warm}}\)}
\KwOut{\(x = \operatorname{prox}_{\tau g + \delta_{\Delta_n}}(v)\); multiplier \(\nu^\star\)}

\BlankLine
\If{\(\tau = 0\)}{
    \Return the Euclidean projection of \(v\) onto \(\Delta_n\) \;
}

\BlankLine
\tcp{Step 1: bracket the simplex multiplier}
Find \([\nu_L,\nu_R]\) such that
\[
\Phi(\nu_L) \ge 1 \ge \Phi(\nu_R),
\qquad
\Phi(\nu) := \mathbf{1}^\top \operatorname{prox}_{\tau g}(v-\nu \mathbf{1}).
\]
Use \(\nu^{\mathrm{warm}}\) as the first guess when available \;

\BlankLine
\tcp{Step 2: initialize the block structure at the right bracket}
Sort \(v\) in descending order \;
Form the shifted scaled vector \((v-\nu_R\mathbf{1})/\tau\) \;
Run the stack-based PAVA routine to obtain the maximal block partition \;
For each block, record its size, penalized-count, current regime, current
fitted value, and current slope \;

\BlankLine
\tcp{Step 3: walk leftward in \(\nu\)}
Set \(\nu \leftarrow \nu_R\) \;
Evaluate the current affine value of \(\Phi(\nu)\) and its slope \;
Create a heap of the next possible events: block regime changes and merges of
adjacent blocks \;

\BlankLine
\While{the root has not yet been identified}{
    \tcp{Between events, \(\Phi\) is affine}
    \If{the affine root of \(\Phi(\nu)=1\) lies before the next queued event}{
        set \(\nu^\star\) to that root and stop \;
    }

    \BlankLine
    \tcp{Otherwise process the next event}
    Pop the next valid event from the heap \;
    Move \(\nu\) to that breakpoint \;
    Update \(\Phi\), its slope, and the affected block data \;
    \If{the event is a regime change}{
        change only that block's regime and slope;
        refresh merge candidates touching that block \;
    }
    \Else{
        merge the two adjacent blocks into one larger block;
        recompute that merged block's regime and slope;
        refresh its neighboring candidates \;
    }
}

\BlankLine
\tcp{Step 4: reconstruct the primal vector}
Evaluate the final block values at \(\nu^\star\) \;
Expand block values back to the sorted coordinates \;
Undo the sorting permutation and return the resulting \(x\) \;
\end{algorithm}

\paragraph*{What is integrated here.}
The solver combines the following ingredients in one implementation:
\begin{itemize}[topsep=0.3em,itemsep=0.2em]
\item an outer proximal augmented Lagrangian loop for \(l \le Ax \le u\);
\item an inner FISTA loop for the primal subproblem;
\item exact evaluation of the free prox \(\operatorname{prox}_{\tau g}\) through
Moreau decomposition;
\item a stack-based PAVA routine for the conjugate-side prox;
\item an exact directional breakpoint walk for the simplex multiplier;
\item warm starts for the multiplier across repeated prox calls;
\item adaptive updates of \(\rho_r\) and \(\eta_r\).
\end{itemize}

\paragraph*{Practical cost of the outer/inner nesting.}
The pseudocode above suggests a full re-solve of \(\textsc{InnerFISTA}\) at
every outer iteration, but the warm-starting built into the implementation
makes the realized cost much lower than that reading implies. Two warm starts
are carried across the nesting:
\begin{itemize}[topsep=0.3em,itemsep=0.2em]
\item the simplex multiplier \(\nu^{\mathrm{warm}}\) is carried both across
successive \textsc{SimplexProx} calls \emph{within} an inner loop and across
outer iterations \(r \to r+1\);
\item the outer iterates \(x^r, p^r, \rho_r, \eta_r\) are carried into the
next \(\textsc{InnerFISTA}\) call as its center and initialization.
\end{itemize}
Once \(\rho_r\) and \(\eta_r\) stop changing rapidly (a few outer rounds in,
by the adaptation rules above), \(y^{(0)} = x^r\) is already close to the next
fixed point, and \(\nu^{\mathrm{warm}}\) already brackets the next root of
\(\Phi\) tightly. In that regime \textsc{InnerFISTA} typically terminates in a
handful of iterations rather than resolving the subproblem from scratch, so
the effective per-outer-iteration cost is close to one gradient step plus one
\textsc{SimplexProx} call, not a full inner solve. The nested-loop structure
is therefore considerably cheaper in practice than the worst-case reading of
the pseudocode suggests.

\paragraph*{Important implementation note.}
The production routine does \emph{not} fall back to bisection if the breakpoint
walk fails. The reference bisection method exists only as a testing oracle. If
the event walk cannot locate the root, the implementation raises an error with
diagnostics instead of silently switching algorithms.

\end{document}
